% ======================================================================
%  udt — plik startowy: SPRAWOZDANIE Z ĆWICZENIA (tryb report)
%
%  Skopiuj ten plik razem z udt-core.sty i udt-report.sty do folderu
%  sprawozdania, zmień nazwę i wypełnij pola. Komentarze opisują każdą
%  funkcję pakietu, z której korzysta sprawozdanie; to, co wspólne
%  z innymi trybami (środowiska, matematyka, warianty), opisano też
%  w plikach streszczenie.tex i skrypt.tex.
%
%  Kompilacja:  pdflatex → biber → pdflatex → pdflatex
%  (biber tylko przy bibliografii; trzeci przebieg rysuje m.in. pionowe
%  kreski dowodów i ustala odsyłacze).
%
%  Wersja rozszerzona z tego samego źródła (zob. WARIANTY niżej):
%      pdflatex -jobname=sprawozdanie-ext sprawozdanie
%
%  Konwencja źródła: jeden akapit = jedna linia, rysunki w rysunki/*.tex.
% ======================================================================

% ----------------------------------------------------------------------
%  KLASA I PAKIET
%  Klasa scrartcl (KOMA-Script); język z opcji klasy (babel): polish albo
%  english. Od języka zależą wszystkie napisy pakietu („Rys.”, „Definicja”,
%  „Oddano: …”, nazwy w \cref) i format liczb siunitx (przecinek dziesiętny).
%
%  Opcje udt-report (można łączyć, np. [chem, full]):
%    plain — środowiska twierdzeń bez kolorowych bloków (wygląd amsthm),
%    chem  — mhchem (\ce) i makra chemiczne \std, \Eo, \redox,
%    full  — włącza wszystkie warianty (zob. WARIANTY).
% ----------------------------------------------------------------------
\documentclass[a4paper, polish, 11pt]{scrartcl}
\usepackage{udt-report}

% Bibliografia (biblatex + biber, numeracja w kolejności cytowania,
% \cite daje [1], [2, 3]). Bez bibliografii usuń tę linię
% i \printbibliography na końcu.
\addbibresource{literatura.bib}

% Własne pakiety dodaje się normalnie, np. schematy elektryczne:
% \usepackage{circuitikz}

% ----------------------------------------------------------------------
%  POLA NAGŁÓWKA
%  Wymagane są tylko \title i \author; puste lub pominięte pole nie
%  zostawia śladu. W trybie report drukowane są:
%    \doctype  — rodzaj dokumentu nad tytułem i w paginie; domyślnie
%                „Sprawozdanie z ćwiczenia”, ale tylko razem z \exercise,
%    \exercise — numer dopisywany do rodzaju dokumentu (dowolny tekst:
%                „2”, „F-10”, „1: część I”),
%    \title    — tytuł, składany wersalikami; dwa wiersze nie nachodzą,
%    \author   — kilku autorów: \author{A \and B}, rozdzieleni kreską,
%    \studentid — numery albumów w tej samej kolejności: {1 \and 2},
%    \supervisor — po „pod nadzorem”, więc w dopełniaczu,
%    \date     — data oddania (domyślnie \today); \date{} usuwa cały
%                wiersz „Oddano: …”,
%    \experimentdate — dopisuje „na podstawie ćwiczenia wykonanego …”.
%  Pola \institution, \course, \group i \place tryb report przyjmuje,
%  ale ich nie drukuje (drukuje je tryb script); warto je wypełnić, bo
%  zostają w źródle i pozwalają przenieść tekst do skryptu.
% ----------------------------------------------------------------------
\institution{Politechnika Gdańska, Wydział Fizyki Technicznej i Matematyki Stosowanej}
\course{Nazwa przedmiotu}
\group{Fizyka Techniczna, Fizyka Stosowana, rok III}
% \doctype{Sprawozdanie z ćwiczenia}   % domyślne, nie trzeba podawać
\exercise{1}
\title{Tytuł ćwiczenia}
\author{Imię Nazwisko \and Imię Nazwisko}
\studentid{000000 \and 000000}
\supervisor{dr. Imię Nazwisko}
\experimentdate{1 października 2026}
\date{\today}

\begin{document}
\maketitle

% ----------------------------------------------------------------------
%  ABSTRAKT I PRZEDMOWA
%  Oba na pełną szerokość, przed \begincolumns, bez nagłówka:
%    abstract — mała, pogrubiona kursywa,
%    preface  — kursywa zamknięta poziomą kreską (np. cel ćwiczenia).
%  Oba są opcjonalne.
% ----------------------------------------------------------------------
\begin{abstract}
Jedno–trzy zdania: co zbadano, jaką metodą i jaki jest główny wynik z niepewnością, np. $g = \qty{9.81(4)}{\metre\per\second\squared}$.
\end{abstract}

% ----------------------------------------------------------------------
%  DWIE KOLUMNY
%  \begincolumns otwiera dwie kolumny w miejscu wywołania; wszystko
%  powyżej zostaje na pełną szerokość. Kolumny zamykają się same na
%  końcu dokumentu. Drugie wywołanie nic nie robi.
% ----------------------------------------------------------------------
\begincolumns

% ----------------------------------------------------------------------
%  ROZDZIAŁY
%  W trybie report rozdziały są nienumerowane; odwołuje się do nich
%  nazwą: \nameref{sec:teoria} → „Zarys teoretyczny”.
%  Równania, rysunki, tabele i środowiska twierdzeń są numerowane
%  1, 2, … w całym dokumencie.
% ----------------------------------------------------------------------
\section{Zarys teoretyczny}\label{sec:teoria}

% Odsyłacze: \cref sam dopisuje nazwę obiektu w odpowiednim języku:
%   \cref{eq:okres} → „równanie (1)”, \cref{rys:schemat} → „Rys. 1”,
%   \cref{tab:pomiary} → „Tab. 1”, \Cref na początku zdania.
% \eqref{eq:okres} daje samo „(1)” — do zdań typu „wzór (1)”.
Dla małych wychyleń okres drgań wahadła matematycznego o długości $l$ wynosi
\begin{equation}
	T = 2\pi\sqrt{\frac{l}{g}}
	\label{eq:okres}
\end{equation}
gdzie $g$ jest przyspieszeniem ziemskim. Z \cref{eq:okres} wynika liniowa zależność $T^2 = (4\pi^2/g)\,l$.

% Przypisy: znacznik z nawiasem („³⁾”), żeby po symbolu nie wyglądał
% jak potęga. Przypisy z wnętrza bloków (twierdzeń, definicji) też
% trafiają na dół strony i mają wspólną numerację.
Wzór~\eqref{eq:okres} jest przybliżeniem małych drgań\footnote{Dla amplitudy $\theta_0$ poprawka pierwszego rzędu wynosi $T \approx T_0\,(1 + \theta_0^2/16)$.}.

% Środowiska twierdzeń: definition, theorem, law, lemma, fact,
% corollary, example, remark (wspólny licznik) oraz proof i aside.
% Szczegóły w streszczenie.tex i skrypt.tex. W sprawozdaniu przydaje się
% głównie law (prawo fizyczne):
\begin{law}[Zasada izochronizmu]\label{prawo:izochronizm}
Okres małych drgań wahadła nie zależy od amplitudy.
\end{law}

\section{Metodologia}\label{sec:metodologia}

% RYSUNKI I TABELE
% W trybie report figure i table stoją dokładnie tam, gdzie je wpisano
% ([H] — w kolumnach nie mogą „pływać”). figure* i table* zajmują obie
% kolumny i trafiają na górę następnej strony.
% Rysunek z pliku: \includegraphics[width=\columnwidth]{rysunki/zdjecie}
% Rysunek TikZ z osobnego pliku: \input{rysunki/schemat}
% Kolory wielkości fizycznych (te same w tekście i w TikZ):
%   udtPos (położenie), udtVel (prędkość), udtAcc (przyspieszenie),
%   udtForce (siła, moment siły), udtMom (pęd, moment pędu),
%   udtAxis (osie, wersory, układy odniesienia).
\begin{figure}
	\centering
	\begin{tikzpicture}[>=Stealth, line cap=round, font=\small]
		\fill[pattern=north east lines, pattern color=udtAxis] (-1,0) rectangle (1,0.2);
		\draw[thick, udtAxis] (-1,0) -- (1,0);
		\draw[dashed, udtAxis] (0,0) -- (0,-2.6);
		\draw[thick] (0,0) -- node[right] {$l$} (-25:2.5) coordinate (m);
		\fill (m) circle (3pt);
		\draw[->, thick, udtForce] (m) -- ++(0,-1) node[right] {$m\va{g}$};
		\draw (0,-0.8) arc[start angle=-90, end angle=-65, radius=0.8];
		\node at (-80:1.1) {$\theta$};
	\end{tikzpicture}
	\caption{Schemat wahadła matematycznego.}
	\label{rys:schemat}
\end{figure}

Długość wahadła $l$ mierzono taśmą mierniczą, a okres $T$ wyznaczano z czasu dziesięciu pełnych drgań (\cref{rys:schemat}).

\section{Wyniki}\label{sec:wyniki}

% Tabela: booktabs (\toprule, \midrule, \bottomrule) i kolumny S
% z siunitx (wyrównanie do przecinka). Nagłówki kolumn S w nawiasach {}.
% Liczby i jednostki: \num{1.23}, \unit{\metre}, \qty{9.81}{\metre\per\second\squared},
% z niepewnością \qty{9.81(4)}{...} → (9,81 ± 0,04) m/s².
\begin{table}
	\centering
	\caption{Okres drgań dla różnych długości wahadła}
	\label{tab:pomiary}
	\begin{tabular}{S[table-format=1.2] S[table-format=1.3] S[table-format=1.3]}
		\toprule
		{$l$ [\unit{\metre}]} & {$T$ [\unit{\second}]} & {$T^2$ [\unit{\second\squared}]} \\
		\midrule
		0.50 & 1.419 & 2.014 \\
		0.75 & 1.738 & 3.021 \\
		1.00 & 2.006 & 4.024 \\
		\bottomrule
	\end{tabular}
\end{table}

% Wykres pgfplots na szerokość kolumny. Rysunek złożony: subfigure,
% podpisy (a), (b) i odsyłacze \cref{rys:wykres-a} → „Rys. 2a”.
% pgfplots nie korzysta z ustawień siunitx: przecinek dziesiętny na osiach
% daje opcja /pgf/number format/use comma.
\begin{figure}
	\centering
	\begin{subfigure}{0.48\columnwidth}
		\centering
		\begin{tikzpicture}
			\begin{axis}[width=\linewidth, xlabel={$l$ [\unit{\metre}]}, ylabel={$T$ [\unit{\second}]}, font=\footnotesize, /pgf/number format/use comma]
				\addplot[only marks, udtPos] coordinates {(0.50,1.419) (0.75,1.738) (1.00,2.006)};
			\end{axis}
		\end{tikzpicture}
		\caption{}
		\label{rys:wykres-a}
	\end{subfigure}\hfill
	\begin{subfigure}{0.48\columnwidth}
		\centering
		\begin{tikzpicture}
			\begin{axis}[width=\linewidth, xlabel={$l$ [\unit{\metre}]}, ylabel={$T^2$ [\unit{\second\squared}]}, font=\footnotesize, /pgf/number format/use comma]
				\addplot[only marks, udtPos] coordinates {(0.50,2.014) (0.75,3.021) (1.00,4.024)};
				\addplot[domain=0.45:1.05, udtForce] {4.025*x};
			\end{axis}
		\end{tikzpicture}
		\caption{}
		\label{rys:wykres-b}
	\end{subfigure}
	\caption{Okres drgań (a) i jego kwadrat (b) w funkcji długości wahadła.}
	\label{rys:wykres}
\end{figure}

Wyniki zebrano w \cref{tab:pomiary}. Zależność $T^2(l)$ na \cref{rys:wykres-b} jest liniowa, zgodnie z~\cref{eq:okres}; z nachylenia prostej $a = 4\pi^2/g$ wyznaczono $g = \qty{9.81(4)}{\metre\per\second\squared}$. Niepewności liczono zgodnie z~\cite{taylor}.

% WARIANTY — treść warunkowa (pełny opis w skrypt.tex).
% Wariant extended jest predefiniowany. Treść w środowisku extended
% pojawia się tylko w wersji rozszerzonej (opcja full albo nazwa
% zadania kończąca się na -ext), oznaczona cienką kreską na marginesie.
% W zwykłej wersji znika, jakby jej nie było. Ukryta treść nie dostaje
% numerów, więc numeracja obu wersji się różni; środowiska z gwiazdką
% (remark*, example*, …) mają ten sam wygląd bez numeru.
\begin{extended}
\begin{remark*}
Pełne obliczenia niepewności: dopasowanie metodą najmniejszych kwadratów daje $a = \qty{4.025(16)}{\second\squared\per\metre}$, skąd $u(g) = g\,u(a)/a$.
\end{remark*}
\end{extended}

\section{Wnioski}\label{sec:wnioski}

Wyznaczona wartość $g$ zgadza się z wartością tablicową w granicach niepewności. Wynik (rozdział \nameref{sec:wyniki}) potwierdza \cref{prawo:izochronizm} w zakresie małych wychyleń.

% Bibliografia: pozycje cytowane przez \cite, w kolejności cytowania.
\printbibliography

\end{document}
